\documentclass[]{article}


% Options for packages loaded elsewhere
\PassOptionsToPackage{unicode}{hyperref}
\PassOptionsToPackage{hyphens}{url}
\PassOptionsToPackage{dvipsnames,svgnames,x11names}{xcolor}

\usepackage{xcolor}
\usepackage[margin=1in]{geometry}
\usepackage{amsmath,amssymb}
\setcounter{secnumdepth}{-\maxdimen} % remove section numbering
\usepackage{iftex}
\ifPDFTeX
  \usepackage[T1]{fontenc}
  \usepackage[utf8]{inputenc}
  \usepackage{textcomp} % provide euro and other symbols
\else % if luatex or xetex
  \usepackage{unicode-math} % this also loads fontspec
  \defaultfontfeatures{Scale=MatchLowercase}
  \defaultfontfeatures[\rmfamily]{Ligatures=TeX,Scale=1}
\fi
\usepackage{lmodern}
\ifPDFTeX\else
  % xetex/luatex font selection
\fi
% Use upquote if available, for straight quotes in verbatim environments
\IfFileExists{upquote.sty}{\usepackage{upquote}}{}
\IfFileExists{microtype.sty}{% use microtype if available
  \usepackage[]{microtype}
  \UseMicrotypeSet[protrusion]{basicmath} % disable protrusion for tt fonts
}{}
\makeatletter
\@ifundefined{KOMAClassName}{% if non-KOMA class
  \IfFileExists{parskip.sty}{%
    \usepackage{parskip}
  }{% else
    \setlength{\parindent}{0pt}
    \setlength{\parskip}{6pt plus 2pt minus 1pt}}
}{% if KOMA class
  \KOMAoptions{parskip=half}}
\makeatother
\usepackage{graphicx}
\makeatletter
\newsavebox\pandoc@box
\newcommand*\pandocbounded[1]{% scales image to fit in text height/width
  \sbox\pandoc@box{#1}%
  \Gscale@div\@tempa{\textheight}{\dimexpr\ht\pandoc@box+\dp\pandoc@box\relax}%
  \Gscale@div\@tempb{\linewidth}{\wd\pandoc@box}%
  \ifdim\@tempb\p@<\@tempa\p@\let\@tempa\@tempb\fi% select the smaller of both
  \ifdim\@tempa\p@<\p@\scalebox{\@tempa}{\usebox\pandoc@box}%
  \else\usebox{\pandoc@box}%
  \fi%
}
% Set default figure placement to htbp
\def\fps@figure{htbp}
\makeatother
\setlength{\emergencystretch}{3em} % prevent overfull lines
\providecommand{\tightlist}{%
  \setlength{\itemsep}{0pt}\setlength{\parskip}{0pt}}
\usepackage{amsmath}
\usepackage{bookmark}
\IfFileExists{xurl.sty}{\usepackage{xurl}}{} % add URL line breaks if available
\urlstyle{same}
\hypersetup{
  pdftitle={ISTx1003 prosjekt: Tittel (gi projektet en informativ tittel)},
  pdfauthor={Gruppenummer},
  colorlinks=true,
  linkcolor={Maroon},
  filecolor={Maroon},
  citecolor={Blue},
  urlcolor={blue},
  pdfcreator={LaTeX via pandoc}}

\title{ISTx1003 prosjekt: Tittel (gi projektet en informativ tittel)}
\author{Gruppenummer}
\date{06 oktober, 2026}

\begin{document}
\maketitle

Rapporten gir maks 50 poeng. Fordeling av poengene vises for hver del.
Hver del blir vurdert etter

\begin{itemize}
\tightlist
\item
  falglig nivå
\item
  korrekthet
\item
  presentasjon og lesbarhet
\end{itemize}

\subsection{Sammendrag/Abstract (max 500 ord) (10
poeng)}\label{sammendragabstract-max-500-ord-10-poeng}

Formålet med sammendraget er å gi en kort og presis oppsummering av
prosjektet deres. Det er en selvstendig tekst som plasseres før selve
rapporten begynner (men skriv den til slutt). Sammendraget bør inneholde
følgende deler:

\begin{enumerate}
\def\labelenumi{\arabic{enumi}.}
\tightlist
\item
  Start med å tydelig beskrive formålet med prosjektet. Hvilket problem
  forsøker dere å løse? Hvilket spørsmål ønsker dere å besvare? Beskriv
  forskningsspørsmålet/hypotesen. Det er viktig å være kortfattet og
  konkret.
\item
  Beskriv dataene dere brukte. Inkluder lenken til datasettet her.
  Hvilket datasett benyttet dere? Deklarer størrelse på datasettet
  (antal observasjoner, antal variabler), typer variabler og
  responsvariabel.
\item
  Beskriv metodene: Hvordan analyserte dere dataene? Hvilke verktøy,
  teknikker eller metoder brukte dere? Regresjon, klassifisering,
  klyngeanalyse? Vær spesifisk, men unngå for mange detaljer.
\item
  Oppsummer de viktigste funnene. I hoveddelen av sammendraget skal dere
  oppsummere de viktigste resultatene fra prosjektet og kort tolke dem
  (altså, hva betyr resultatene?). Fremhev de mest betydningsfulle
  funnene, og gi nok informasjon til at leseren får et inntrykk av hva
  dere har oppdaget.
\item
  (Valgfritt) Fremhev betydningen av resultatene. Forklar hvorfor
  og/eller hvordan funnene dine er viktige. Fremhev eventuelle nye eller
  uventede funn, og forklar hvordan de bidrar til økt forståelse av
  temaet.
\end{enumerate}

\subsection{Introduksjon (10 poeng)}\label{introduksjon-10-poeng}

\begin{itemize}
\item
  Gi en kort introduksjon til den overordnede problemstillingen eller
  oppgaven dere har valgt, samt datasettet dere bruker. Dette kan for
  eksempel være en klassifiseringsoppgave (f.eks. å forutsi om en
  pasient utvikler hjertesykdom eller ikke) eller en prediksjonsoppgave
  (f.eks. å forutsi prisen på en Airbnb). Definer tydelig avgrensningen
  for prosjektet deres. Hvilket spesifikt problem forsøker dere å løse?
\item
  Beskriv kilden til datasettet, og oppgi en referanse til hvor
  datasettet kommer fra.
\item
  Beskriv formålet med prosjektet mer detaljert. Hvilke konkrete
  spørsmål ønsker dere å besvare i prosjektet? Prøver dere å finne
  metoden som gir best resultater, eller en metode som gir gode
  resultater samtidig som den er enkel og effektiv å bruke? Hvem er
  målgruppen? Forsøker dere å undersøke sammenhenger mellom ulike
  variabler? Prøver dere å finne viktige prediktorer for
  klassifiseringen? Eller ønsker dere å oppnå en bedre og mer
  innsiktsfull forståelse av et bestemt tema eller fagområde?
\end{itemize}

\subsection{Beskrivende dataanalyse (5
poeng)}\label{beskrivende-dataanalyse-5-poeng}

Her bør du gjennomføre en kort deskriptiv dataanalyse for å få en
oversikt over datasettet. Det betyr at man lager noen visualiseringer
for å forstå dataene bedre, og at man beregner noen enkle ting. Fokuser
på det som er relevant for modelleringen, og bruk sunn fornuft. For mye
detaljer eller figurer uten forklaring og tydelige akseetiketter er ikke
nyttige for leseren.

Du kan for eksempel:

\begin{itemize}
\tightlist
\item
  Presentere deskriptiv statistikk, som gjennomsnitt, median,
  standardavvik og varians, for å beskrive datasettets sentraltendens,
  variasjon og fordeling.
\item
  Bruke spredningsdiagrammer (scatter plots) og korrelasjonsmatriser for
  å undersøke sammenhenger mellom ulike variabler.
\item
  Bruke histogrammer for å vise fordelingen til relevante variabler.
\item
  Bruke boksplott (box plots) for å undersøke fordeling, variasjon og
  eventuelle uteliggere.
\end{itemize}

Resultatene bør forklares i teksten, med særlig fokus på funn som kan
være relevante for den videre modelleringen. Skriv en kort legende under
hver grafikk, slik at det blir tydelig hva grafikken viser.

\subsection{Metoder (10 poeng)}\label{metoder-10-poeng}

\begin{itemize}
\item
  Beskriv metodene dere bruker i prosjektet, og forklar i detalj hvordan
  dere har anvendt dem. Dere bør bruke metoder fra mist 2 av de 3
  temaene i kurset (regresjon, klassifikasjon, klyngeanalyse).
\item
  Forklar kort hvorfor metodene passer for datasett og problemstillingen
  dere har valgt.
\item
  \textbf{Det er viktig at dere bare bruker metoder som er tilpasset
  nivået i dette kurset -- bruk bare metoder dere lærer om her. Hold det
  enkelt og realistisk, ikke bruk komplekse metoder som dere ikke
  forstår, bare fordi KI foreslår det.}
\item
  Beskriv hvordan dere evaluerer metodene, hvis det er mulig. Evaluer
  modellen(e) ved å bruke passende verktøy, for eksempel residualanalyse
  eller vurdering av prediksjonsytelse (som sensitivitet og spesifisitet
  i en binær klassifiseringsoppgave).
\item
  (Valgfritt) Diskuter mulige begrensninger ved metodene og de valgte
  evalueringsmålene. Vurder for eksempel om et bestemt mål kan gi et
  misvisende bilde av modellens ytelse, eller om det finnes egenskaper
  ved datasettet som gjør en metode mindre egnede.
\end{itemize}

\subsection{Resultater (10 poeng)}\label{resultater-10-poeng}

\begin{enumerate}
\def\labelenumi{\arabic{enumi}.}
\item
  Presenter resultatene på en tydelig og strukturert måte. Dette kan for
  eksempel gjøres ved hjelp av tabeller, grafer eller andre
  visualiseringer som gjør funnene lettere å forstå.
\item
  Tolk resultatene. Tolkningen bør tilpasses regresjons-,
  klassifiserings- eller klyngeanalyse-oppgaven du har valgt.
\item
  Diskuter begrensninger og forbehold som er viktige å ta hensyn til når
  resultatene skal tolkes. Dette kan for eksempel være begrensninger ved
  datasettet, modellene eller evalueringsmetodene.
\item
  (Valgfritt) Se på mulige videre analyser. Beskriv eventuelt
  alternative eller forbedrede metoder som kunne vært brukt for å
  analysere dataene i fremtiden, og hvordan disse potensielt kunne gitt
  ytterligere innsikt.
\end{enumerate}

\subsection{Diskusjon og konklusjon (5
poeng)}\label{diskusjon-og-konklusjon-5-poeng}

Oppsummer de viktigste funnene fra prosjektet deres. Hva fant du ut, og
hvilke var de viktigste innsiktene du fikk gjennom analysen? Fokuser på
de mest relevante resultatene fra modelleringen og analysen.

\newpage

\subsection{Deklarasjon om KI-hjelpemidler (maks 1 side -- teller ikke
til de 10 sidene i
rapporten)}\label{deklarasjon-om-ki-hjelpemidler-maks-1-side-teller-ikke-til-de-10-sidene-i-rapporten}

Besvar folgende sprørsmål:

1.) Har det i utarbeidingen av denne rapporten blitt anvendt KI-baserte
hjelpemidler? Hvis ja, beskriv hvilke verktøy dere brukte (Copilot, NTNU
GPT,\ldots).

2.) Er en eller flere av ideene eller fremgangsmåtene i oppgaven
foreslått av KI? Hvis ja, beskriv her:

3.) Tekst: Har du brukt KI for

\begin{enumerate}
\def\labelenumi{\alph{enumi})}
\tightlist
\item
  Stavekontroll?
\item
  Språkvask?
\item
  Tekstgenerering?
\item
  Skriveassistanse.
\end{enumerate}

Hvis du svarer ja på minst en av a)-d), spesifiser bruken her:

4.) Kode og algoritmer: Programmeringsassistanse. Er deler av
koden/algoritmene som i) fremtrer direkte i rapporten eller ii) har
blitt anvendt for produksjon av resultater slik som figurer, tabeller
eller tallverdier blitt generert av KI? Hvis ja, spesifiser bruken her:

5.) Bilder og figurer, bildegenerering. Er ett eller flere av
bildene/figurene i rapporten blitt generert av KI? Hvis ja, beskriv
bruken her:

6.) Beskrive alle andre aspekter av arbeidet som KI ble brukt til.
Beskriv her:

7.) Ved innlevering av prosjektet deklarerer vi at vi er kjent med NTNUs
regelverk: Det er ikke tillatt å generere besvarelser eller tekst ved
hjelp av KI og levere den helt eller delvis som eget arbeid. Vi har
derfor redegjort for all anvendelse av KI enten direkte i rapporten,
eller her. Svar her:

\newpage

\subsection{Referanser og kode (teller ikke til de 10 sidene i
rapporten)}\label{referanser-og-kode-teller-ikke-til-de-10-sidene-i-rapporten}

Her kan dere legge til alle referanser og kode som dere har brukt for å
generere resultatene i rapporten.

\end{document}
